\chapter*{Preface}

These notes for the \textbf{Computing Methods for Experimental Physics and Data Analysis} course for Master's Degree students at the University of Pisa are based on the in-person classes given during the 2025-26 academic year. The \textit{Advanced Machine Learning} module (3 cfu) is included, but you won't find any chapters about the \textit{Medical Physics} or \textit{Particle Physics} ones. Images, graphs and tables are mostly taken from the course slides, available on the E-Learning page. Notebooks and code snippets were also part of the course material. 

I tried to create a document that was as comprehensive and as student-friendly as possible while still maintaining a certain rigor, especially in the most "mathematical" parts. Being a student myself, I am no expert in the field of Machine Learning or parallel programming. As of now these notes have not been reviewed by any of the professors of the course, so tread lightly. There will be errors of every possible kind: spelling errors, grammatical and syntactical errors, conceptual errors and so on. If (or better, \textit{when}) you find one, please report it in the Issues page of the GitHub repository where I store and publish these notes (\url{https://github.com/LorenzoNepi/notes_CMEPDA}) or correct it yourself by opening a Pull Request. 

As you can see from the table of contents, each topic covered corresponds to a whole chapter and each section contains (more or less) a single lesson. I preferred to almost always follow the chronological order of the course. One exception to this is Sec. \ref{sec:GNN_intro}, which contains both the introductory and advanced lessons on the topic of GNNs.

Here is the full process that gave birth to this document:
\begin{itemize}
    \item[-] During the year, I collected my own notes, the corresponding presentation slides and the audio transcripts automatically generated by Teams (if the lesson was recorded).
    \item[-] I passed all the material to Gemini Pro so that it could generate the skeleton of the notes with very precise and simple formatting instructions.
    \item[-] I copied all the Notebooks and ran them on my Colab, adding comments and Markdown cells to make the code as readable as possible. 
    \item[-] I personally reviewed, expanded and rewrote every section starting from the Gemini output.
    \item[-] I passed the sections one by one to Gemini one more time to catch the biggest errors and mistakes (don't worry, there will be others).   
\end{itemize}
If you decide to trust me and my notes at your own peril, I have only one thing left to say:\\ \textit{May the force be with you}.